\scanpage{10}

% archive sheet number: 9

Supposons maintenant que $C$ soit lui aussi un $A$-module libre (de sorte que la condition relative aux quasi-primitifs sera une conséquence de l'autre condition). Je dis que \uline{pour que} \deleted{$\varphi$}

\[
\varphi:\deleted{\unsure{B}}C\otimes_A C\longrightarrow\operatorname{Hom}_A(U,C)\tag*{\unsure{$(*)$|$(\times)$}}
\]

\uline{soit un isomorphisme}, \uline{il faut et il suffit que l'hom analogue défini par} $\varphi$

\[
\psi:C\otimes U\longrightarrow\operatorname{Hom}(C,C)\tag*{$(**)$}
\]

\uline{par la condition}

\[
\boxed{\psi(x\otimes u).y=x\,u(y)}
\]

\margin{?}

\uline{soit un isomorphisme}. [Donc, si $B$ et $C$ sont tous deux libres de type fini, on a deux caractérisations \deleted{\unsure{des fibrés principaux de groupe}} \inserted{équivalentes} \inserted{du fait que $C$ soit} \deleted{fibré} un fibré principal homogène sous $C$]

Supposons en effet que $\varphi$ soit un isomorphisme, i.e. \deleted{étendons les scalaires à $A'=C$, alors}

\[
\deleted{C^{A'}\otimes_{A'}C^{A'}\longrightarrow\operatorname{Hom}_{A'}(U^{A'},C^{A'})}
\]

\deleted{est encore un \unsure{$B$|$C$|$B^{A'}$} est ``fibré hom. principal sous $B=U'$'' dans la catégorie des $A$-algèbres comm. Faisant le produit \deleted{fibré} par $A'=C$, on trouve que \inserted{ds la catégorie des $A'$-algèbres} le fibré $C^{A'}=C\otimes_A A'$ sur $B^{A'}=U'\otimes_A A'=(U^{A'})'$ est isomorphe à $B^{A'}$ lui-même, [car on a pris $A'=C$].}
